\section*{4. What Survived, What Stayed As Exploration}

\noindent
\begin{minipage}[t]{0.48\linewidth}
\softpanel{tealaccent}{%
\textbf{What clearly survived}
\begin{itemize}
    \item strong tabular feature artifacts, especially the target-stats substrate
    \item careful branch combination instead of single-model heroics
    \item lighter neural or meta layers used as finishers, not as total replacements
    \item splitting work sensibly across Mac/MPS and HPC/CUDA
    \item treating Kaggle submissions as calibration points rather than blind search
\end{itemize}}
\end{minipage}\hfill
\begin{minipage}[t]{0.48\linewidth}
\softpanel{ink}{%
\textbf{What mostly stayed as exploration}
\begin{itemize}
    \item fully replacing the strong tabular core with pure classifier-only neural ideas
    \item the most aggressive feature factories when they added more complexity than signal
    \item tree-teacher and distillation-style branches as final leaders
    \item the professor classical screen as a direct replacement for the best existing lanes
    \item any branch whose exact history was not preserved well enough to defend with artifacts
\end{itemize}}
\end{minipage}

\vspace{0.8em}

Overall, the project improved most when it emphasized disciplined iteration rather than isolated novelty. The strongest gains usually came from better feature artifacts, cleaner pruning, stronger tabular branches, and careful combination of those branches.

The professor-inspired lane still mattered. It widened the search space, improved the refreshed target-stats branch slightly, and forced the project to think more carefully about missingness and feature selection. But it behaved more like a good upstream research lane than like a full new winner.

For the final release, the overnight family was used as the main published pipeline rather than the absolute best local blend. Its performance remained strong, while its branch structure and artifact trail were substantially more stable than those of the later experimental hybrids.

\vspace{0.5em}

\softpanel{tealaccent}{%
\textbf{Final takeaways}
\begin{enumerate}
    \item The score story was mostly a story about better artifacts and better combinations, not about abandoning trees.
    \item The best local branch did exceed the released overnight lane, but the overnight family remained the more stable and reproducible final pipeline.
    \item Some old exact results are gone, but the overall trend is still clear from the saved score tables and logs.
    \item If the project were continued, the most sensible next step would be to keep the target-stats backbone, keep the professor lane as an upstream research branch, and only retune the final combiner when the upstream source truly improves.
\end{enumerate}}
